%! TeX program = lualatex
% Two pictures about how to space practice, shared by the reading-week slides
% and the lesson of the same week.
%
%   page 1 -- how OFTEN a technique is practised.  Ten problems either way; on
%             top they are split among five shortcuts (two repetitions each),
%             below they all use one systematic approach (ten repetitions).
%   page 2 -- in what ORDER techniques are practised.  Twelve problems and
%             three techniques either way; blocked on top, interleaved below.
%
% One tile is one practice problem and its colour is the technique used, so the
% two pictures share a vocabulary and read as a pair.  Both rows of a picture
% hold the same number of tiles at the same positions on purpose: the only
% thing that changes between the rows is the quantity the picture is about, so
% neither can be misread as "do more problems".
\documentclass[tikz]{standalone}
\input{../tikz.tex.preamble}
\input{../typography.tex.preamble}

\tikzset{
  practice tiles/.style={
    tile/.style={rounded corners=2pt, draw=WesternGrey,
                 minimum width=0.75cm, minimum height=0.6cm, inner sep=0pt},
    caption/.style={anchor=west, font=\small},
  },
}

\begin{document}
% ---------------------------------------------------------------------
% Page 1 -- repetition count.  Five shortcuts twice each against one
% systematic approach ten times.
\begin{tikzpicture}[practice tiles]
  \node[caption] at (-0.375, 1.75) {Five shortcuts, \(2\) repetitions each};
  \foreach \c [count=\n from 0] in {%
      WesternSky, WesternSky, WesternSpring, WesternSpring, WesternVivid,
      WesternVivid, WesternTiger, WesternTiger, WesternOrchid, WesternOrchid}
    \node[tile, fill=\c!70] at ({\n*0.85}, 1.1) {};

  \node[caption] at (-0.375, 0.15) {One systematic approach, \(10\) repetitions};
  \foreach \n in {0,...,9}
    \node[tile, fill=main!40] at ({\n*0.85}, -0.5) {};
\end{tikzpicture}
% ---------------------------------------------------------------------
% Page 2 -- ordering.  Blocked practice on top, interleaved below; the second
% ordering is the one that forces a choice of technique before every problem.
\begin{tikzpicture}[practice tiles]
  \node[caption] at (-0.375, 1.75) {Learning a technique: one type of problem at a time};
  \foreach \c [count=\n from 0] in {%
      WesternSky, WesternSky, WesternSky, WesternSky,
      WesternSpring, WesternSpring, WesternSpring, WesternSpring,
      WesternTiger, WesternTiger, WesternTiger, WesternTiger}
    \node[tile, fill=\c!70] at ({\n*0.85}, 1.1) {};

  \node[caption] at (-0.375, 0.15) {Closer to the exam: mixed together};
  \foreach \c [count=\n from 0] in {%
      WesternSky, WesternTiger, WesternSpring, WesternTiger,
      WesternSky, WesternSpring, WesternTiger, WesternSky,
      WesternSpring, WesternSky, WesternTiger, WesternSpring}
    \node[tile, fill=\c!70] at ({\n*0.85}, -0.5) {};
\end{tikzpicture}
\end{document}
